% !TEX root = ../../tree_notes.tex
% =====================================================================
%  2.3  Bottom-Up Recursion: returning values up          sec:c02-bottomup
%  Exports: def:c02-bottomup  prop:c02-aggregate  lst:c02-bottomup-tmpl
%           prob:c02-tilt  prob:c02-prune  prob:c02-univalue
% =====================================================================
\section{Bottom-Up Recursion: Returning Values Up}\label{sec:c02-bottomup}

\begin{topicintro}
In a \term{bottom-up} recursion each node asks its children for an answer, combines the
two, and \emph{returns} the result to its own parent -- the return value is a subtree
aggregate. The height and size recurrences of \thmref{c01-height-recursion} were already
bottom-up; here we generalise them and add the one idiom that unlocks the hardest tree
mediums: a function that \emph{returns} one quantity (so the parent can use it) while
\emph{silently updating a global best} (the real answer). That is exactly how diameter,
tilt and longest-path problems collapse into a single DFS.

\textbf{Prerequisites.} \secref{c02-model} (structural recursion) and
\thmref{c01-height-recursion} (return-and-combine).

\textbf{After this section you can}
\begin{itemize}[nosep]
  \item tell a bottom-up problem instantly: the answer at a node is a function of its two
        subtrees alone;
  \item return a subtree aggregate (sum, height, pruned pointer, chain length) and combine
        it correctly;
  \item use the "return one value, update a second global answer" split for
        diameter/tilt/longest-path;
  \item solve binary tree tilt, prune tree, and longest univalue path in one pass each.
\end{itemize}
\end{topicintro}

\usedin{Chapter 9 (diameter and centroid), Chapter 10 (tree DP) and Chapter 20 (advanced
tree DP)}{tree DP \emph{is} bottom-up recursion with a richer returned state; mastering the
return-plus-global split now is literally the first lesson of every DP-on-trees chapter.}

\subsection{Intuition}

\begin{analogybox}[An org chart reporting up]
A VP does not decide the company total herself: she waits for both directors under her to
report their numbers, adds them to her own, and passes one total up to the CEO. The value
travels from the bottom (leaves) to the top (root) -- the opposite direction to
\secref{c02-topdown}. Sometimes, though, a manager also shouts "biggest team so far!" to a
whiteboard on the wall: that side note is the global-best update made during the upward pass.
\end{analogybox}

Take subtree sums. To know the sum at a node you must already know the sums of both
children, so the information is forced to flow upward. We draw the returns on the tree.

\begin{figure}[H]
\centering
\begin{tikzpicture}[diagram]
  \node[troot] (r) at (0,0) {$1$};
  \node[tnode] (l) at (-1.8,-1.4) {$2$};
  \node[tnode] (c) at (1.8,-1.4) {$3$};
  \node[tleaf] (a) at (-2.6,-2.8) {$4$};
  \node[tleaf] (b) at (-1.0,-2.8) {$5$};
  \node[tleaf] (d) at (1.8,-2.8) {$6$};
  \draw[tedge] (r) -- (l); \draw[tedge] (r) -- (c);
  \draw[tedge] (l) -- (a); \draw[tedge] (l) -- (b); \draw[tedge] (c) -- (d);
  \node[annot, text=successgreen] at (-2.9,-3.3) {$4$};
  \node[annot, text=successgreen] at (-1.0,-3.3) {$5$};
  \node[annot, text=successgreen] at (-2.3,-1.9) {$11$};
  \node[annot, text=successgreen] at (1.8,-3.3) {$6$};
  \node[annot, text=successgreen] at (2.4,-1.9) {$9$};
  \node[annot, text=successgreen] at (0.2,0.7) {$21$};
\end{tikzpicture}
\caption{Subtree sums returned upward (green): leaves return their own value, node $2$
returns $2{+}4{+}5=11$, node $3$ returns $3{+}6=9$, the root returns $1{+}11{+}9=21$. The
parent cannot be computed before the children -- the definition of bottom-up.}\label{fig:c02-bottomup}
\end{figure}

\subsection{Definition and the Aggregate Property}

\begin{definition}{Bottom-up recursion}{c02-bottomup}
A recursion whose \emph{return value} is a function of the return values of the children
plus the node itself: $g(v)=\phi\big(v,\ g(v_{\text{left}}),\ g(v_{\text{right}})\big)$
with a base value at \nil. The answer to the whole problem is either $g(\text{root})$ or a
quantity recorded while computing the $g$'s.
\end{definition}

\begin{proposition}{Bottom-up well-definedness}{c02-aggregate}
Because a tree has no cycles (\defref{c01-tree}) every recursion path ends at a \nil{}
base case, so the combination $\phi$ is applied over a finite, acyclic call graph and $g(v)$
depends only on the strict subtree $\subtree(v)$ -- the aggregate is always defined.
\end{proposition}
\begin{proofbox}
$\subtree(v)$ is finite and, excluding $v$, consists of the strict subtrees of $v$'s
children, which are strictly smaller. Induction on $\size(v)$ (\thmref{c01-height-recursion}
is exactly this for $\height$): the children's $g$ are defined by the hypothesis (smaller
subtrees), and $\phi$ combines finitely many of them with $v$. The \nil{} base terminates
the recursion.
\end{proofbox}

\subsection{The Template and the Global-Best Split}

\begin{lstlisting}[caption={Bottom-up skeleton: return the aggregate the parent needs, and separately update the answer that lives "through" the current node}, label={lst:c02-bottomup-tmpl}]
// g(r) = the aggregate a parent can reuse (height, sum, chain length, pruned pointer).
// best = the real answer, which may pass THROUGH r and so cannot be a single-child return.
Aggregate dfs(TreeNode* r, Best& best) {
    if (r == nullptr) return IDENTITY;        // 0 for sums, 0/-1 for heights, ...
    Aggregate L = dfs(r->left,  best);
    Aggregate R = dfs(r->right, best);
    best = std::max(best, combine(L, R, r));  // answer using BOTH sides at r
    return fold(L, R, r);                    // single aggregate sent UP
}
\end{lstlisting}

\begin{notebox}[Why the split is necessary]
A recursion returns \emph{one} value, but the best path/tilt/answer at $r$ may use both
children at once (a path bending through $r$). If you tried to return that, the parent would
double-count. So you return the extendable part (one side) and record the through-part in a
global -- the exact move behind diameter (\chref{c09}) and \probref{c02-univalue} below.
\end{notebox}

\begin{complexitybox}
One DFS, $\bigO{n}$ time, $\bigO{h}$ stack. The whole point of bottom-up is that each
child's aggregate is computed once and reused by the parent, so an $\bigO{n}$-per-node
naive recomputation (like the $O(n^2)$ balanced-check warned about in \probref{c01-balanced})
drops to $\bigO{n}$.
\end{complexitybox}

\subsection{Worked Examples}

\begin{example}{Reading returns off \figref{c02-bottomup}}{c02-bu-ex}
The returned sums are $4,5$ at $2$'s children, so node $2$ returns $11$; $3$ returns $9$;
the root returns $21$. If instead we wanted the max subtree sum, we would update a global
with each return and read $\max(4,5,11,6,9,21)=21$.
\end{example}

\begin{tracebox}
\code{dfs} for tilt visits $4,5$ then $2$: at $2$ the tilt is $|5-4|=1$, returns sum $11$;
at $3$ tilt $|6-0|=6$, returns $9$; at the root tilt $|9-11|=2$. Total tilt $1+6+2=9$.
\end{tracebox}

\subsection{Common Pitfalls}

\begin{warningbox}[Returning the answer instead of the aggregate]
For path-through-node problems, returning $\text{combine}(L,R,r)$ (the through value)
poisoning the parent is the most common error. Return only what a parent can extend -- the
better of the two children -- and keep the through value in the global.
\end{warningbox}

\begin{warningbox}[Wrong identity for the null base case]
Sums return $0$, node-heights return $0$, edge-heights return $-1$, "contains a $1$" returns
\code{nullptr}/\code{false}. Using $-1$ where $0$ is the identity makes every subtree off by
one -- the same off-by-one flagged in \secref{c01-anatomy}.
\end{warningbox}

\subsection{Solved Problems}

\solvedproblem{c02-tilt}{Binary Tree Tilt}{LeetCode 563}
  {https://leetcode.com/problems/binary-tree-tilt/}{LC Easy}
  {bottom-up subtree sum, global accumulator}

\begin{problemstatement}[Binary Tree Tilt]
The tilt of a node is $|\text{sum}(\text{left})-\text{sum}(\text{right})|$. Return the sum
of the tilts of all nodes.
\end{problemstatement}

\solstep{Understanding the Problem}
Tilt needs each subtree's \emph{sum}, which is a clean bottom-up aggregate; the total tilt
is a value living at each node, so we return the sum and add the tilt into a global --
exactly the split in \coderef{c02-bottomup-tmpl}.

\solstep{Implementation}
\begin{lstlisting}[caption={Binary Tree Tilt (LeetCode 563)}, label={lst:c02-tilt}]
class Solution {
    int total = 0;
    int sum(TreeNode* r) {
        if (!r) return 0;                       // null contributes sum 0
        int L = sum(r->left), R = sum(r->right);
        total += abs(L - R);                    // tilt recorded on the way up
        return r->val + L + R;                  // return the subtree sum
    }
public:
    int findTilt(TreeNode* root) { sum(root); return total; }
};
\end{lstlisting}

\solstep{Complexity}
\begin{complexitybox}
Time $\bigO{n}$, memory $\bigO{h}$. Note the two distinct roles: \code{sum} returns the
subtree total, \code{total} accumulates the answer -- folding them into one return would
lose the per-node tilt.
\end{complexitybox}

\solvedproblem{c02-prune}{Binary Tree Pruning}{LeetCode 814}
  {https://leetcode.com/problems/binary-tree-pruning/}{LC Medium}
  {bottom-up decision, return subtree pointer}

\begin{problemstatement}[Binary Tree Pruning]
The tree holds only $0$ and $1$. Remove every subtree that does not contain a $1$, and
return the root of the pruned tree.
\end{problemstatement}

\solstep{Understanding the Problem}
Whether to keep a node depends only on its subtree: keep $r$ iff $r$\texttt{->val}$=1$ or
one of its (already-pruned) children survived. That is a bottom-up predicate, so we recurse
first, then decide, returning \code{nullptr} to delete a whole branch.

\solstep{Key Observations}
\begin{observation}
After pruning the children, $r$ must go iff $\text{val}=0$ \emph{and} both children are now
\nil.
\end{observation}
\begin{observation}
A $1$ always survives even if both children are pruned away.
\end{observation}

\solstep{Implementation}
\begin{lstlisting}[caption={Binary Tree Pruning (LeetCode 814)}, label={lst:c02-prune}]
class Solution {
public:
    TreeNode* pruneTree(TreeNode* r) {
        if (!r) return nullptr;
        r->left  = pruneTree(r->left);    // prune below first
        r->right = pruneTree(r->right);
        if (r->val == 0 && !r->left && !r->right) return nullptr;  // dead branch
        return r;
    }
};
\end{lstlisting}

\solstep{Complexity}
\begin{complexitybox}
Time $\bigO{n}$, memory $\bigO{h}$. The returned pointer doubles as the aggregate (survives
or not) and the rebuilt child link -- a bottom-up recursion that mutates the tree in place.
\end{complexitybox}

\solvedproblem{c02-univalue}{Longest Univalue Path}{LeetCode 687}
  {https://leetcode.com/problems/longest-univalue-path/}{LC Easy}
  {bottom-up chain length, global best, path-through-node}

\begin{problemstatement}[Longest Univalue Path]
Return the number of \emph{edges} on the longest path whose nodes all share one value. The
path need not pass the root and may bend at any node.
\end{problemstatement}

\solstep{Understanding the Problem}
Mirror the diameter idea (\probref{c01-balanced}'s cousin in \secref{c01-anatomy}): \code{dfs}
returns the longest univalue \emph{chain} going straight down from $r$ (extendable by $r$'s
parent), while the global answer takes the best path \emph{through} $r$, which can join the
left and right chains when they match $r$\texttt{->val}. This is the through-vs-return split
of \coderef{c02-bottomup-tmpl} verbatim.

\solstep{Implementation}
\begin{lstlisting}[caption={Longest Univalue Path (LeetCode 687)}, label={lst:c02-univalue}]
class Solution {
    int best = 0;
    int chain(TreeNode* r) {
        if (!r) return 0;
        int L = chain(r->left), R = chain(r->right);
        int leftExt = 0, rightExt = 0;
        if (r->left  && r->left->val  == r->val) leftExt  = L + 1;   // extendable chains
        if (r->right && r->right->val == r->val) rightExt = R + 1;
        best = std::max(best, leftExt + rightExt);                   // path through r
        return std::max(leftExt, rightExt);                          // return one side up
    }
public:
    int longestUnivaluePath(TreeNode* root) { chain(root); return best; }
};
\end{lstlisting}

\solstep{Complexity}
\begin{complexitybox}
Time $\bigO{n}$, memory $\bigO{h}$. Returning $\max(\text{leftExt},\text{rightExt})$ but
scoring $\text{leftExt}+\text{rightExt}$ in the global is the whole trick -- returning the
sum would let the parent count a path that bends twice.
\end{complexitybox}

\solstep{Pitfalls and Variants}
\begin{warningbox}[Counting edges, not nodes]
The answer is edges; a chain of $k$ matching nodes scores $k-1$. Adding $1$ only when we
\emph{extend} to a child (as above) keeps the count in edges automatically.
\end{warningbox}
Variant: \emph{Diameter of Binary Tree} is the same code with lengths in place of value
matching; the general tree version is \chref{c09}.

\subsection{Practice Problems}

\begin{practiceset}
  \practice{Distribute Coins in Binary Tree}{LeetCode 979}
    {https://leetcode.com/problems/distribute-coins-in-binary-tree/}{LC Medium}
    {bottom-up: return the net surplus/deficit a subtree must move through its parent edge; add its absolute value to the answer.}
  \practice{House Robber III}{LeetCode 337}
    {https://leetcode.com/problems/house-robber-iii/}{LC Medium}
    {return a pair (rob this node, skip this node) per subtree; combine the two children's pairs.}
  \practice{Delete Nodes And Return Forest}{LeetCode 1110}
    {https://leetcode.com/problems/delete-nodes-and-return-forest/}{LC Medium}
    {post-order; when a child is deleted, its subtree roots become new trees -- collect them.}
  \practicehh{Path Sum III}{LeetCode 437}
    {https://leetcode.com/problems/path-sum-iii/}{LC Medium}
    {count downward paths (any start/end) summing to target; brute force is two DFSs.}
    {prefix sums on the root-to-node path: carry a hash map of running-sum counts down, and for each node add how many earlier prefixes equal current$-$target. A top-down map plus a bottom-up count.}
\end{practiceset}
